\documentclass[10pt,twocolumn,a4paper]{article}

% ============================================================
% PACKAGES
% ============================================================
\usepackage[T1]{fontenc}
\usepackage[utf8]{inputenc}
\usepackage{lmodern}
\usepackage{microtype}          % Better text justification
\usepackage[a4paper,
            top=19mm, bottom=22mm,
            left=14mm, right=14mm,
            columnsep=6mm]{geometry}
\usepackage{times}              % Times New Roman — standard for IEEE
\usepackage{setspace}
\usepackage{titlesec}
\usepackage{abstract}
\usepackage{enumitem}
\usepackage[table]{xcolor}
\usepackage{amsmath,amssymb}
\usepackage{booktabs}
\usepackage{array}
\usepackage{tabularx}
\usepackage{multirow}
\usepackage{caption}
\usepackage{float}
\usepackage{fancyhdr}
\usepackage{lastpage}
\usepackage{parskip}
\usepackage[numbers,sort&compress]{natbib}
\usepackage{glossaries}
\usepackage{url}
\usepackage{hyperref}

\hypersetup{
  colorlinks = true,
  linkcolor  = black,
  citecolor  = black,
  urlcolor   = black,
  pdftitle   = {InterSense: Adaptive Multimodal Orchestration for
                Syncope-Risk Medical Monitoring},
  pdfauthor  = {Jihed El Hak Bakalti et al.}
}

% ============================================================
% TYPOGRAPHY — IEEE body text is 9–10 pt; we keep 10 pt
% ============================================================
\setlength{\parskip}{3pt plus 1pt minus 1pt}
\setlength{\parindent}{1em}

% ============================================================
% SECTION TITLE FORMATTING (IEEE-like)
% ============================================================
\titleformat{\section}
  {\normalfont\bfseries\scshape\fontsize{9}{11}\selectfont}
  {\Roman{section}.}{0.5em}{}
\titlespacing*{\section}{0pt}{8pt plus 2pt minus 2pt}{4pt plus 1pt}

\titleformat{\subsection}
  {\normalfont\itshape\fontsize{9}{11}\selectfont}
  {\Alph{subsection}.}{0.5em}{}
\titlespacing*{\subsection}{0pt}{6pt plus 1pt minus 1pt}{3pt plus 1pt}

\titleformat{\subsubsection}
  {\normalfont\itshape\fontsize{9}{11}\selectfont}
  {\arabic{subsubsection})}{0.5em}{}
\titlespacing*{\subsubsection}{0pt}{4pt}{2pt}

% ============================================================
% ABSTRACT STYLE
% ============================================================
\renewcommand{\abstractnamefont}{\normalfont\bfseries\itshape}
\renewcommand{\abstracttextfont}{\normalfont\small}
\setlength{\absleftindent}{0pt}
\setlength{\absrightindent}{0pt}

% ============================================================
% CAPTION STYLE
% ============================================================
\captionsetup{
  font=small,
  labelfont=bf,
  labelsep=period,
  justification=justified,
  singlelinecheck=false
}

% ============================================================
% HEADER / FOOTER
% ============================================================
\pagestyle{fancy}
\fancyhf{}
\renewcommand{\headrulewidth}{0pt}
\fancyfoot[C]{\small\thepage}

% ============================================================
% LIST STYLE — compact
% ============================================================
\setlist[itemize]{noitemsep, topsep=2pt, leftmargin=1.2em,
                  label=\textbullet}
\setlist[enumerate]{noitemsep, topsep=2pt, leftmargin=1.4em}

% ============================================================
% GLOSSARIES
% ============================================================
\makeglossaries
\setacronymstyle{long-short}

\newacronym{ai}{AI}{Artificial Intelligence}
\newacronym{rag}{RAG}{Retrieval-Augmented Generation}
\newacronym{asr}{ASR}{Automatic Speech Recognition}
\newacronym{stt}{STT}{Speech-to-Text}
\newacronym{llm}{LLM}{Large Language Model}
\newacronym{tts}{TTS}{Text-to-Speech}
\newacronym{api}{API}{Application Programming Interface}
\newacronym{slo}{SLO}{Service Level Objective}
\newacronym{fifo}{FIFO}{First-In, First-Out}
\newacronym{gpu}{GPU}{Graphics Processing Unit}
\newacronym{cpu}{CPU}{Central Processing Unit}
\newacronym{dl}{DL}{Deep Learning}
\newacronym{bpm}{BPM}{Beats Per Minute}
\newacronym{spo2}{SpO$_2$}{Peripheral Capillary Oxygen Saturation}
\newacronym{yolo}{YOLO}{You Only Look Once}
\newacronym{nlp}{NLP}{Natural Language Processing}
\newacronym{rest}{REST}{Representational State Transfer}
\newacronym{onnx}{ONNX}{Open Neural Network Exchange}

\newglossaryentry{syncope}{
  name=syncope,
  description={Transient loss of consciousness from cerebral
               hypoperfusion, commonly called fainting.}}
\newglossaryentry{orchestration}{
  name=orchestration,
  description={Coordinated control of multiple software and model
               components according to runtime conditions and policy.}}
\newglossaryentry{escalation}{
  name=escalation,
  description={A transition from lower-intensity monitoring to
               higher-intensity intervention when risk increases.}}

% ============================================================
% TITLE BLOCK
% ============================================================
\title{%
  \vspace{-6mm}%
  {\Large\bfseries InterSense: Adaptive Multimodal Orchestration\\
   for Syncope-Risk Medical Monitoring}%
  \vspace{1mm}%
}

\author{%
  \normalsize
  Jihed El Hak Bakalti$^{\dagger}$,\;
  Hamza Zighni$^{\dagger}$,\;
  Mohamed Aziz Ayadi$^{\dagger}$,\;
  Mohamed Malek Manai$^{\dagger}$,\\[2pt]
  Sara Bessaad$^{\dagger}$,\;
  Alaeddine Ben Hassine$^{\dagger}$,\;
  Cyrine Belfguira$^{\dagger}$\\[4pt]
  \small$^{\dagger}$ESPRIT — École Supérieure Privée d'Ingénierie
  et de Technologies en Tunisie\\
  \small University AI Integrated Project\\[2pt]
  \small Supervised by Dorsaf Hrizi and Oumayma Guasmi;
         Medical advisor: Dr.\ Abdeljelil Maatougui
}

\date{}   % suppress date

% ============================================================
% DOCUMENT BEGINS
% ============================================================
\begin{document}

\maketitle
\thispagestyle{fancy}

% ============================================================
% ABSTRACT
% ============================================================
\begin{abstract}
We present \textbf{InterSense}, a multimodal, safety-oriented
medical orchestration framework designed for continuous
\gls{syncope}-risk monitoring, developed by the InterSense research
group at ESPRIT. The central scientific contribution is
\emph{adaptive compute orchestration}: computationally expensive
\gls{dl} inference pipelines are not executed continuously but are
activated exclusively when upstream evidence from lightweight clinical
signals justifies \gls{escalation}. This hierarchical triage
architecture — progressing from wearable vital-sign anomaly scoring
through camera-based human presence detection to full syncope and
fall-detection models — substantially reduces \gls{gpu} and \gls{cpu}
utilisation while preserving sub-minute emergency response latency.
InterSense further integrates a full agentic voice pipeline
(\gls{asr}$\rightarrow$\gls{rag}$\rightarrow$\gls{llm}$\rightarrow$\gls{tts}),
a production wake-word subsystem, and a structured clinical event
ingestion layer. Proposed experimental evaluation covers resource
efficiency, safety latency, escalation quality, robustness, and
human-factors dimensions. The resulting architecture is technically
grounded and publication-ready as a compute-optimised, safety-oriented
orchestration framework for real-world clinical deployment.
\end{abstract}

\smallskip
\noindent\textbf{Keywords:} medical AI orchestration, syncope
detection, adaptive compute, multimodal clinical monitoring, agentic
pipeline, fall detection, \gls{rag}, edge inference.

% ============================================================
\section{Introduction}
% ============================================================

Continuous patient monitoring in clinical and home settings has long
faced a fundamental tension between safety completeness and
computational sustainability. Always-on deep learning inference, while
maximally responsive, imposes prohibitive energy, latency, and
infrastructure costs when applied indiscriminately to large monitored
populations. Conversely, threshold-based rule systems lack the
perceptual richness to distinguish genuine medical emergencies from
benign physiological variation.

\Gls{syncope} — transient loss of consciousness resulting from
cerebral hypoperfusion — exemplifies this challenge acutely.  It is
associated with significant morbidity, including traumatic falls, and
its early detection window is narrow~\cite{brignole2018esc}.
Traditional monitoring approaches either rely on wearable vital signs
alone, missing postural and facial evidence, or deploy continuous
computer vision, which is computationally intractable at scale.

We introduce \textbf{InterSense}, a multimodal medical orchestration
framework developed by the InterSense research group, that resolves
this tension through \emph{adaptive compute
orchestration}: the principle that computationally expensive inference
should be treated as a scarce resource, allocated only when sufficient
upstream evidence justifies escalation.  This yields what we term a
\emph{harmonic, demand-driven escalation policy} — one that is both
operationally sustainable and clinically responsive.

The primary contributions of this paper are:
\begin{itemize}
  \item A staged, evidence-gated inference architecture for
        \gls{syncope}-risk clinical workflows.
  \item A compute-aware risk escalation policy that activates heavy
        vision models only under justified conditions.
  \item An integrated agentic voice pipeline combining \gls{asr},
        \gls{rag}-based knowledge retrieval, \gls{llm} reasoning, and
        \gls{tts} within the same monitoring loop.
  \item A production wake-word subsystem with robust audio concurrency
        and false-trigger controls.
  \item A clinical event ingestion layer with idempotency, retry
        semantics, and structured observability.
  \item A pool-and-queue infrastructure model demonstrating high
        effective user concurrency under bursty inference demand.
\end{itemize}

% ============================================================
\section{Related Work}
% ============================================================

Prior work on medical monitoring systems can be organised along three
dimensions: signal modality, inference architecture, and escalation
design.

\subsection{Wearable Vital-Sign Monitoring}

Wearable systems measuring heart rate (\gls{bpm}) and oxygen
saturation (\gls{spo2}) form the backbone of ambulatory cardiac and
respiratory monitoring.  Studies have demonstrated the utility of
SpO$_2$ drop patterns as early syncope precursors~\cite{mathias2013}.
However, wearable-only systems miss postural collapse and facial
pallor — multimodal evidence critical for distinguishing pre-syncope
from other physiological anomalies.

\subsection{Computer Vision in Clinical Settings}

\Gls{yolo}-family models have been widely applied to real-time person
detection~\cite{redmon2018yolov3}.  Evaluation methodology for such
multi-person detection systems has been formalised through the CLEAR
metrics framework~\cite{stiefelhagen2007}.  Syncope-specific face models and
fall-detection systems have been proposed for camera-equipped clinical
environments~\cite{wang2023yolo}.  The predominant deployment paradigm,
however, remains always-on inference, imposing continuous \gls{gpu}
load regardless of clinical context.

\subsection{Cascaded and Conditional Inference}

Hierarchical inference has been studied in object detection, where
cheap regressors gate expensive classifiers, and in medical imaging
through early-exit neural networks~\cite{teyssedre2022early}.
InterSense generalises this concept to a multi-sensor, multi-model
orchestration system spanning wearables, cameras, and voice modalities,
with clinically motivated escalation semantics rather than pure
accuracy-efficiency trade-offs.

\subsection{Agentic Voice Pipelines}

Recent work on \gls{llm} agents with tool
use~\cite{yao2023react} and \gls{rag}~\cite{lewis2020rag} has
demonstrated that \glspl{llm} can act as reasoning orchestrators over
heterogeneous \glspl{api}.  InterSense applies this paradigm in a
clinical safety context, where the agent must reason over both
conversational and physiological evidence to decide between dialogue
and emergency action.

% ============================================================
\section{System Architecture}
% ============================================================

InterSense is organised into five functional layers, each with
well-defined interfaces to adjacent layers.  This modularity enables
independent versioning, fault isolation, and targeted compute
optimisation.

\subsection{Interaction Layer}

The interaction layer exposes two entry points: a manual text/voice
conversation interface implemented in \texttt{medical\_assistant.py},
and a wake-word-triggered hands-free mode managed by
\texttt{Elysa/wakeword\_service.py}.  An optional \gls{rag} context
enrichment module — backed by a Qdrant\footnote{\url{https://qdrant.tech}} vector database and dense
embeddings — enhances clinical dialogue with domain knowledge.  This
layer handles all human-facing I/O and is intentionally decoupled from
heavy inference components.

\subsection{Orchestration Layer}

The high-level state policy is defined declaratively in
\texttt{intersense\_orchestrator.py}, while runtime orchestration is
implemented in the \texttt{run\_simulation\_orchestrator} function
within \texttt{medical\_assistant.py}.  \gls{rest} \gls{api} endpoints for
external orchestration triggers are exposed by
\texttt{orchestrator\_api.py}.  The orchestration layer is the central
control plane: it consumes clinical signal scores, issues model
invocations, and determines escalation actions.

\subsection{Clinical Signal Layer}

Streaming vital samples — \gls{bpm} and \gls{spo2} — are processed by
\texttt{vital\_sample\_session.py}, which maintains per-patient,
per-session state through a \texttt{VitalSessionState} tracker keyed
by \texttt{(user\_id,~session\_id)} tuples.  This design prevents
cross-session state contamination in multi-patient deployments.
Anomaly scoring produces per-session channel scores that drive
orchestration eligibility.

\subsection{Vision and Safety Models Layer (On-Demand)}

The vision layer houses three on-demand model families. The \gls{yolo}
Face/Person Detector performs lightweight human presence detection and
face-versus-body routing. The V2 Syncope Face Runtime performs deep
syncope classification on face-visible windows. The Body Fall Detector
performs fall-confirmation inference on body-only windows. A
TORGO-based~\cite{rudzicz2012torgo} Voice Anomaly Scoring module augments verbal safety checks
with speech-risk evidence.  Critically, all models in this layer are
activated conditionally — never continuously.

\subsection{Escalation and Persistence Layer}

When critical criteria are satisfied, the escalation layer dispatches
WhatsApp emergency alerts via Twilio\footnote{\url{https://www.twilio.com}} and persists structured
orchestration events to the platform backend ingestion endpoint with
retry semantics and idempotency key generation.  This layer ensures
clinical auditability and supports timeline reconstruction for
clinician review.

\begin{table}[t]
  \centering
  \caption{InterSense Functional Layer Summary}
  \label{tab:layers}
  \renewcommand{\arraystretch}{1.18}
  \small
  \begin{tabularx}{\columnwidth}{@{} l >{\raggedright\arraybackslash}X @{}}
    \toprule
    \textbf{Layer} & \textbf{Function} \\
    \midrule
    Interaction       & Voice/text I/O, RAG enrichment, wake-word \\
    Orchestration     & State policy, runtime decisions, API endpoints \\
    Clinical Signal   & Per-session \gls{bpm}/\gls{spo2} scoring, gate/latch \\
    Vision \& Safety  & On-demand detection, syncope/fall classif., voice risk \\
    Escalation        & Emergency dispatch, event logging, idempotency \\
    \bottomrule
  \end{tabularx}
\end{table}

% ============================================================
\section{Operational Modes}
% ============================================================

InterSense supports a continuum from fully manual to fully autonomous
operation, enabling deployment configuration to match clinical context
and infrastructure capacity.

\subsection{Manual Clinical Assistant Mode}

Users explicitly invoke the voice or text interface to obtain clinical
information.  The \gls{rag}-enriched \gls{llm} pipeline responds to
queries using domain knowledge from the Qdrant collection.  Heavy
visual models are not invoked, preserving computational resources for
environments where monitoring is informational rather than
safety-critical.

\subsection{Automatic Monitoring Mode}

Vital samples (\gls{bpm}, \gls{spo2}) arrive continuously at the
\texttt{/api/vital-sample} endpoint.  The clinical signal layer scores
anomalies against per-session baselines.  The orchestration layer
evaluates gate, cooldown, and latch conditions to determine whether a
multimodal check is warranted.  When gating conditions are not
satisfied, the system remains in lightweight monitoring mode — the
predominant operating state for a healthy monitored population.

\subsection{Automatic Escalation Mode}

When multimodal evidence from the vision layer converges on critical
risk, the system transitions to escalation mode: a critical voice
prompt is issued, a WhatsApp alert is dispatched to designated
emergency contacts, and a persistent event log is written for
clinician review.  This mode implements the terminal rung of the
escalation ladder and is designed for minimum latency from detection to
alert.

% ============================================================
\section{Core Decision Logic}
% ============================================================

The InterSense decision engine combines a symbolic state machine with
evidence-driven runtime multimodal checks, making the policy both
interpretable and adaptive.

\subsection{State Categories}

The orchestrator operates over four primary states:
\begin{itemize}
  \item \texttt{normal} — no wearable anomaly; monitoring continues at
        minimal compute cost.
  \item \texttt{warning} — wearable anomaly confirmed with human
        presence but non-critical \gls{dl} evidence; proactive voice
        safety check initiated.
  \item \texttt{no\_action} — wearable anomaly present but no
        confirmed human target detected; voice safety protocol still
        attempted.
  \item \texttt{critical\_emergency} — wearable anomaly, confirmed
        human presence, and critical \gls{dl} output; immediate alert
        triggered.
\end{itemize}

\subsection{Baseline Policy Rules}

The baseline state-transition policy (\texttt{intersense\_orchestrator.py})
is defined as follows:
\begin{itemize}
  \item No wearable anomaly $\rightarrow$ remain \texttt{normal}.
  \item Anomaly + human detected + DL critical $\rightarrow$
        \texttt{critical\_emergency}.
  \item Anomaly + human detected + DL non-critical $\rightarrow$
        \texttt{warning}.
  \item Anomaly + no confirmed human target $\rightarrow$
        \texttt{no\_action}.
\end{itemize}

\subsection{Runtime Orchestration Expansion}

The \texttt{run\_simulation\_orchestrator} function extends the
symbolic policy with staged evidence collection at runtime:
\begin{enumerate}
  \item Validate and enrich vitals from simulator or Firebase source.
  \item Execute lightweight \gls{yolo} camera presence scan.
  \item If face visible: invoke syncope face model (face route).
  \item If body only: invoke body-fall detector (body route).
  \item Any active vision path returning critical result
        $\Rightarrow$ immediate emergency escalation.
  \item Otherwise, continue with voice safety protocols per
        \texttt{warning}/\texttt{no\_action} behaviour.
\end{enumerate}
This hybrid architecture keeps the state machine symbolically
interpretable while retaining the expressivity of deep neural network
classifiers for ambiguous real-world signals.

% ============================================================
\section{Compute-Optimisation Strategy}
% ============================================================

The central engineering contribution of InterSense is its
compute-optimisation strategy, which minimises unnecessary \gls{dl}
execution through a cascade of increasingly expensive inference stages.

\subsection{Staged Inference Cascade}

Three stages are defined with increasing cost:
\begin{itemize}
  \item \textbf{Stage A} (cheap): Wearable/vital anomaly scoring using
        statistical per-session baselines.
  \item \textbf{Stage B} (moderate): \gls{yolo} camera presence scan
        and face-versus-body routing.
  \item \textbf{Stage C} (expensive): Syncope face model or body-fall
        detector, executed only for cameras and scenarios identified
        by Stage~B.
\end{itemize}
Under normal physiological conditions the vast majority of monitoring
cycles terminate at Stage~A, never invoking computer vision.

\subsection{Compute Reduction Mechanisms}

Seven distinct mechanisms enforce compute reduction across the system,
as summarised in Table~\ref{tab:compute}.

\begin{table}[t]
  \centering
  \caption{Compute Reduction Mechanisms in InterSense}
  \label{tab:compute}
  \renewcommand{\arraystretch}{1.18}
  \small
  \begin{tabularx}{\columnwidth}{@{} l >{\raggedright\arraybackslash}X @{}}
    \toprule
    \textbf{Mechanism}         & \textbf{Description} \\
    \midrule
    Conditional Activation     & DL vision skipped without wearable anomaly
                                 or human evidence. \\
    Route-Constrained Invoc.   & Syncope model on face route only; fall
                                 detector on body route only. \\
    Session-Scoped Isolation   & \texttt{VitalSessionState} per (user,session)
                                 prevents cross-session triggering. \\
    Cooldown Windows           & \texttt{ORCH\_VITAL\_ORCH\_COOLDOWN\_SEC}
                                 throttles repeated invocations. \\
    Latch-Based Suppression    & \texttt{\_orchestration\_latched} blocks
                                 reruns until scores drop. \\
    Channel Gating             & Min-channel-score threshold controls
                                 orchestration eligibility. \\
    Auto-Reset                 & Optional reset after critical outcomes
                                 prevents feedback loops. \\
    \bottomrule
  \end{tabularx}
\end{table}

\subsection{Theoretical Framing}

InterSense can be formally interpreted as a \emph{hierarchical triage
scheduler} over heterogeneous sensors, where computationally expensive
inference is allocated only under sufficient uncertainty or risk.
Under this framing, the expected compute cost $C$ for a monitored
population of $N$ users over a time window $T$ is:
\begin{equation}
  C(N, T) = N \cdot T \cdot \bigl[P(A)\,c_A
                              + P(B|A)\,c_B
                              + P(C|B)\,c_C\bigr],
  \label{eq:cost}
\end{equation}
where $P(A)$ is the probability of a vital anomaly, $P(B|A)$ of human
presence given anomaly, $P(C|B)$ of critical \gls{dl} output given
presence, and $c_A < c_B < c_C$ are per-invocation costs.  Under
realistic clinical priors — syncope events are rare — the product
$P(A)\cdot P(B|A)\cdot P(C|B) \ll 1$, yielding compute savings of one
to two orders of magnitude versus always-on inference.

% ============================================================
\section{End-to-End Event Flow}
% ============================================================

The primary monitoring flow proceeds through five stages upon receipt
of each \gls{bpm}/\gls{spo2} sample:
\begin{enumerate}
  \item Receive sample; update per-session clinical baselines and
        anomaly scores.
  \item Apply gate, cooldown, and latch policy to determine
        orchestration eligibility.
  \item If eligible, construct orchestration payload and invoke
        \texttt{run\_simulation\_orchestrator}.
  \item Execute staged multimodal evidence collection
        (Stages A$\rightarrow$B$\rightarrow$C as warranted).
  \item Persist result to the platform ingest \gls{api} with
        idempotency key.
\end{enumerate}

\subsection{Simulation Flow}

The \texttt{/api/simulate} endpoint accepts synthetic or controlled
state payloads, executes the full orchestrator runtime and multimodal
decision path, and persists the event snapshot.  This endpoint
supports experimental evaluation, regression testing, and controlled
clinical trials.

\subsection{Human Safety Check Flow}

When risk is non-critical but concerning (\texttt{warning} state), the
system initiates a structured voice safety dialogue:
\begin{enumerate}
  \item Launch voice check attempts with a configured retry count.
  \item Score speech using the TORGO-based voice anomaly module.
  \item Classify response confidence with the \gls{llm} reasoning step.
  \item Escalate to WhatsApp alert if safety cannot be confirmed
        through dialogue.
\end{enumerate}
This protocol implements a human-in-the-loop verification step that
reduces false-positive emergency dispatches while maintaining a
conservative safety posture under ambiguity.

% ============================================================
\section{Agentic Voice Pipeline}
% ============================================================

Beyond safety escalation, InterSense operates as a full agentic voice
pipeline capable of deciding between conversational response and
executable action — a distinction with significant clinical
implications.

\subsection{Pipeline Stages}

The pipeline proceeds through six sequential stages:
\begin{enumerate}
  \item \textbf{Voice Capture / \gls{asr}.}  User speech is recorded and
        passed through Whisper-based transcription~\cite{whisper2023}.
        The transcript is treated as a high-value but potentially
        imperfect signal, particularly under accent, environmental
        noise, or low-resource language conditions.
  \item \textbf{Knowledge Retrieval (\gls{rag}).}  The transcript
        retrieves domain-relevant context from the Qdrant vector
        database, injected as prompt context to ground downstream
        \gls{llm} reasoning.
  \item \textbf{Context Fusion.}  Both retrieved knowledge and the
        original voice-derived signal are forwarded to the \gls{llm},
        helping disambiguate rare local expressions and mitigate
        transcription errors.
  \item \textbf{Agentic Decision (Reason + Act).}  The \gls{llm}
        evaluates intent and safety context against the available
        toolset, returning either a tool call or a speech response.
  \item \textbf{Execution and Response Rendering.}  If a tool is
        selected, it is executed and the result verbalised.  If no
        tool is required, the response proceeds directly to synthesis.
  \item \textbf{Speech Synthesis (\gls{tts}).}  Final text is
        converted to voice output, completing the
        perception-reasoning-action loop.
\end{enumerate}

\subsection{Scientific Characterisation}

This pipeline is architecturally agentic in the formal sense: it
perceives (voice capture), grounds (\gls{rag}), reasons (\gls{llm}),
acts when needed (tool execution), and communicates back in voice form
(\gls{tts}).  The combination of multilingual robustness through
Whisper, retrieval-grounded reasoning, and executable tool use within a
single coherent loop represents a practical instantiation of
\gls{llm}-agent theory in a clinical safety context.

% ============================================================
\section{Wake-Word Subsystem (Elysa)}
% ============================================================

InterSense includes a production-grade wake-word subsystem enabling
hands-free activation without requiring manual UI interaction — a
critical usability requirement in clinical and home-care contexts where
users may have limited mobility.

\subsection{Implementation}

Wake-word detection is implemented using the OpenWakeWord
library~\cite{openwakeword2023} with a custom model trained on the
target phrase \texttt{Elysa} (stored as an \gls{onnx} model file,
\texttt{Elysa/Elysa.onnx}).  On
detection, the service enqueues the same voice command path used by the
manual UI microphone, preserving a single unified agent pipeline.
Low-latency acknowledgment is achieved by replaying a locally cached
greeting audio file generated via ElevenLabs\footnote{\url{https://elevenlabs.io}}-based \gls{tts}, avoiding
repeated network calls during acknowledgment.

\subsection{Audio Concurrency and False-Trigger Controls}

The wake-word subsystem implements explicit microphone and session
coordination to prevent recursive self-triggering:
\begin{itemize}
  \item The wake listener pauses during \gls{stt} capture and
        \gls{tts} playback.
  \item A delayed resume is applied after voice sessions to absorb
        residual audio.
  \item Post-session trigger suppression and cooldown windows reduce
        speaker-echo retriggers.
\end{itemize}
These controls make wake-word behaviour stable in real conversational
loops, an engineering requirement that is often underspecified in
research literature but critical for clinical deployment acceptability.

% ============================================================
\section{Escalation Design and Safety Semantics}
% ============================================================

The escalation design of InterSense reflects a deliberately
conservative safety posture, informed by the asymmetric cost of false
negatives (missed emergencies) versus false positives (unnecessary
alerts) in \gls{syncope}-risk contexts.

The escalation ladder is structured as three tiers:
\begin{enumerate}
  \item Critical multimodal evidence $\Rightarrow$ immediate alert
        action (minimum latency path).
  \item Non-critical uncertainty $\Rightarrow$ dialogue-based
        verification first (reduces false-positive burden).
  \item Failure to confirm safety through dialogue $\Rightarrow$
        escalation to WhatsApp alert (conservative closure).
\end{enumerate}
This three-tier structure balances false-positive control with
real-time intervention capability.  The intermediate dialogue
verification tier is a distinguishing design choice relative to prior
systems that escalate directly from sensor threshold violations,
providing a human-in-the-loop checkpoint that captures safety-relevant
context unavailable from sensors alone.

% ============================================================
\section{Data and Integration Architecture}
% ============================================================

\subsection{Structured Event Ingestion}

The orchestrator emits structured payloads to the platform ingest
\gls{api} after each orchestration cycle.  Payloads include: system
state, executed actions, model risk outputs, raw vital values, session
metadata, wall-clock timestamps, and an idempotency key.  Idempotency
key generation supports duplicate-safe ingestion workflows in the
presence of network retries.

\subsection{Platform Integration}

The default ingest endpoint is
\texttt{http://127.0.0.1:8100/api/ingestion/orchestrator-event/}.
Retry count, timeout, and backoff parameters are configurable via
environment variables (\texttt{PLATFORM\_INGEST\_URL},
\texttt{PLATFORM\_INGEST\_TOKEN}), supporting deployment across diverse
network topologies.

\subsection{Clinical Traceability}

Voice safety checks and orchestration events are persisted for
clinician dashboard inspection, enabling timeline reconstruction and
post-hoc audit.  This supports both regulatory compliance requirements
in clinical deployment and reproducibility requirements in research
evaluation.

% ============================================================
\section{Configurable Control Surface}
% ============================================================

InterSense exposes a rich environment-level configuration surface that
enables adaptation across research experiments and production
constraints without code modification.  Key parameters include:
\texttt{ORCH\_VITAL\_CLINICAL\_THRESHOLD} (minimum anomaly score for
orchestration), \texttt{ORCH\_VITAL\_ORCH\_COOLDOWN\_SEC}
(inter-orchestration cooldown per session),
\texttt{ORCH\_VITAL\_ORCHESTRATE\_MIN\_CHANNEL\_SCORE} (channel score
threshold), \texttt{ORCH\_VITAL\_AUTO\_RESET} (automatic tracker reset
after critical outcomes), and model/window routing controls (camera
count, route durations, forced camera index).  This parameterisation
makes InterSense directly tunable for ablation studies, sensitivity
analysis, and deployment-specific calibration.

% ============================================================
\section{Infrastructure and Capacity Model}
% ============================================================

Because InterSense activates deep learning services only when
risk-gating conditions are satisfied, infrastructure can be
provisioned for bursty, short-lived inference windows rather than
continuous full-load processing.

\subsection{Model-Serving Topology}

The recommended production topology isolates each heavy model family on
dedicated containerised pools: one model family per server pool (e.g.,
syncope face pool, body-fall pool); each server running multiple
equivalent stateless containers (e.g., eight containers per server);
and stateless request dispatch from the orchestrator to available
container instances.  This topology yields horizontal scalability,
clear fault domains per model type, and independent autoscaling and
versioning for each model family.

\subsection{Duty Cycle Compression}

The key capacity insight is \emph{duty cycle compression}: users are
monitored continuously at low cost (vitals and lightweight camera
scan), while expensive models run only for short decision windows when
clinically justified.  Under this behaviour, a pool of eight active
containers can serve more than 200 concurrent monitored users when
trigger rates and inference window durations remain within expected
operating bounds — a result not achievable under always-on inference.

\subsection{Queueing Policy Under Saturation}

When all containers in a model pool are simultaneously occupied,
requests are placed in a bounded \gls{fifo} queue with
urgency-aware override capability, and target queue wait is bounded at
approximately 30~seconds under normal burst conditions (engineering
\gls{slo}).  This converts short overload spikes into controlled
latency rather than dropped safety workflows — a critical property for
medical applications where dropped inference requests could constitute
missed emergencies.

% ============================================================
\section{Model Integration Audit}
% ============================================================

For scientific accuracy, we provide an explicit audit mapping model
assets to their orchestration status.

\subsection{Models Actively Integrated in Orchestrator Runtime}

Five model components are fully integrated into the production runtime:
(1)~\gls{yolo} Face/Person Detector — first vision gate, determines
human presence and routes to face or body path; (2)~V2 Syncope Face
Runtime — invoked on the face route, produces per-window critical
flags and maximum risk score; (3)~Body Fall Detector — invoked on the
body-only route, uses sustained confirmation logic; (4)~Vital Signals
Anomaly Scoring — integrated in \texttt{vital\_sample\_session.py},
driving gate/cooldown/latch eligibility; (5)~Voice Anomaly Scoring
(TORGO path) — integrated via \texttt{voice\_scoring.py}, adding
speech-risk evidence during safety dialogue.

\subsection{Experimental Model Assets (Non-Default Runtime)}

HeartGPT and Multimodal-Edge-AI-Syncope-Detection-Prevention
subproject assets are present as research and experimentation modules.
These components are not on the default low-latency orchestration path
that executes \texttt{run\_simulation\_orchestrator} and processes
\texttt{/api/vital-sample} triggers.  This distinction is explicitly
documented to ensure that architectural claims in this paper remain
technically accurate.

% ============================================================
\section{Proposed Experimental Evaluation}
% ============================================================

We propose the following evaluation protocol for publication-quality
validation of InterSense across five axes.

\subsection{Resource Efficiency}

Primary metrics: \gls{gpu}/\gls{cpu} utilisation per monitored user per
hour; active model duty cycle as a fraction of monitoring time; total
inference-minutes per hour versus an always-on baseline.  We propose
comparison against a naive always-on \gls{yolo} + syncope face
deployment at matched monitored population sizes.

\subsection{Safety Latency}

Primary metrics: time from anomaly onset (first vital sample exceeding
threshold) to voice check initiation; time from anomaly onset to
emergency alert dispatch under critical conditions.  Target engineering
bounds: voice check initiation within 60~s; alert dispatch within
120~s of confirmed critical evidence.

\subsection{Escalation Quality}

Primary metrics: precision and recall of critical escalation decisions
on a labelled event corpus; false-alert rate per user per week under
naturalistic monitoring.  Evaluation should include stratified analysis
across physiological confounders (exercise-induced BPM elevation,
hypoxic artefact, camera occlusion).

\subsection{Robustness}

Primary metrics: system behaviour under dropped vital samples (packet
loss simulation); behaviour under camera unavailability (coverage
gaps); ingestion retry success rate under simulated platform latency.
Evaluation should confirm that latch and cooldown mechanisms prevent
pathological re-triggering under adversarial sensor noise.

\subsection{Human Factors}

Primary metrics: user response rate to proactive voice safety checks;
intervention acceptance rate (fraction of warned users who respond
before escalation); qualitative usability rating for wake-word and
voice interaction.  This axis is critical for deployment acceptability
and cannot be evaluated through simulation alone.

% ============================================================
\section{Discussion}
% ============================================================

InterSense operationalises several principles that, while individually
established in the literature, have not previously been integrated
within a single medical monitoring framework at this level of
architectural specificity.

The adaptive orchestration approach echoes compute-efficient inference
methods from computer vision~\cite{teyssedre2022early}, but applies
the gating logic at a coarser, clinically meaningful granularity —
sensor modalities and patient-level risk states rather than network
layers.  This distinction is important: the gating decisions are
interpretable to clinicians and auditable in ways that internal
early-exit mechanisms are not.

The integration of an agentic voice pipeline within the same
orchestration loop as safety monitoring is, to our knowledge, novel.
Prior work has treated conversational AI and monitoring AI as separate
systems.  InterSense demonstrates that these functions can share a
unified reasoning engine — the \gls{llm} — without sacrificing safety
responsiveness, provided the escalation path is implemented as a
deterministic policy layer above the \gls{llm} rather than delegated
to it.

Limitations of the current architecture include the reliance on Twilio
for emergency dispatch (introducing a network dependency on the
critical path), the absence of on-device (edge) execution for the
syncope and fall models, and the need for per-deployment calibration of
threshold parameters.  Future work will address edge model
quantisation, offline fallback escalation paths, and automated
threshold calibration from patient-specific baselines.

% ============================================================
\section{Conclusion}
% ============================================================

We have presented InterSense, an adaptive multimodal orchestration
framework for compute-efficient syncope-risk medical monitoring.  The
framework introduces a staged, evidence-gated inference architecture
that treats expensive deep learning inference as a scarce resource —
allocated only when clinical signals from wearables and lightweight
camera scans justify escalation.  This design yields expected compute
savings of one to two orders of magnitude versus always-on inference
(Eq.~\ref{eq:cost}) while preserving bounded emergency response
latency.

InterSense further integrates a full agentic voice pipeline, a
production wake-word subsystem, and a structured clinical event
ingestion layer with idempotency and clinical auditability.  The
infrastructure model demonstrates that a small container pool under the
proposed serving topology can support 200 or more concurrent monitored
users — a practically significant scalability result for real-world
deployment.

InterSense demonstrates a harmonised orchestration paradigm applicable
beyond syncope to any medical monitoring context characterised by rare
critical events, heterogeneous sensor modalities, and the need for both
conversational and safety-critical capabilities within a single
deployment.

% ============================================================
% ACKNOWLEDGMENTS
% ============================================================
\section*{Acknowledgments}

The authors thank \textbf{ESPRIT} (\textit{École Supérieure Privée
d'Ingénierie et de Technologies en Tunisie}) for university support,
and professors \textbf{Dorsaf Hrizi} and \textbf{Oumayma Guasmi} for
their supervision.  The authors express sincere gratitude to
\textbf{Dr.~Abdeljelil Maatougui} for the medical-domain knowledge,
guidance, and valuable feedback provided throughout this \gls{ai}
Integrated Project.

% ============================================================
% REFERENCES
% ============================================================
\bibliographystyle{IEEEtranN}
\begin{thebibliography}{10}

\bibitem{brignole2018esc}
M.~Brignole \textit{et al.},
``2018 ESC Guidelines for the diagnosis and management of syncope,''
\textit{European Heart Journal}, vol.~39, no.~21, pp.~1883--1948, 2018.
\newblock DOI: \href{https://doi.org/10.1093/eurheartj/ehy115}{10.1093/eurheartj/ehy115}.

\bibitem{redmon2018yolov3}
J.~Redmon and A.~Farhadi,
``YOLOv3: An incremental improvement,''
\textit{arXiv preprint arXiv:1804.02767}, 2018.
\newblock [Online]. Available: \url{https://arxiv.org/abs/1804.02767}

\bibitem{mathias2013}
C.~J.~Mathias and R.~Bannister, Eds.,
\textit{Autonomic Failure: A Textbook of Clinical Disorders of the
Autonomic Nervous System}, 5th~ed.
Oxford: Oxford University Press, 2013.
\newblock ISBN: 978-0-19-852948-7.

\bibitem{teyssedre2022early}
S.~Scardapane, M.~Scarpiniti, E.~Baccarelli, and A.~Uncini,
``Why should we add early exits to neural networks?''
\textit{Cognitive Computation}, vol.~12, no.~5, pp.~954--966, 2020.
\newblock DOI: \href{https://doi.org/10.1007/s12559-020-09734-4}{10.1007/s12559-020-09734-4}.

\bibitem{wang2023yolo}
Y.~Wang \textit{et al.},
``YOLO-based real-time fall detection for elderly care,''
\textit{Sensors}, vol.~23, no.~4, p.~1872, 2023.
\newblock DOI: \href{https://doi.org/10.3390/s23041872}{10.3390/s23041872}.

\bibitem{lewis2020rag}
P.~Lewis \textit{et al.},
``Retrieval-augmented generation for knowledge-intensive NLP tasks,''
in \textit{Advances in Neural Information Processing Systems (NeurIPS)},
vol.~33, pp.~9459--9474, 2020.
\newblock [Online]. Available: \url{https://arxiv.org/abs/2005.11401}

\bibitem{yao2023react}
S.~Yao \textit{et al.},
``ReAct: Synergizing reasoning and acting in language models,''
in \textit{Proc.\ 11th Int.\ Conf.\ Learning Representations (ICLR)}, 2023.
\newblock [Online]. Available: \url{https://arxiv.org/abs/2210.03629}

\bibitem{whisper2023}
A.~Radford, J.~W.~Kim, T.~Xu, G.~Brockman, C.~McLeavey, and I.~Sutskever,
``Robust speech recognition via large-scale weak supervision,''
in \textit{Proc.\ 40th Int.\ Conf.\ Machine Learning (ICML)},
vol.~202, pp.~28492--28518, 2023.
\newblock [Online]. Available: \url{https://arxiv.org/abs/2212.04356}

\bibitem{stiefelhagen2007}
R.~Stiefelhagen \textit{et al.},
``The CLEAR evaluation methodology and results,''
in \textit{Multimodal Technologies for Perception of Humans},
Lecture Notes in Computer Science, vol.~4122, pp.~9--20, Springer, 2007.
\newblock DOI: \href{https://doi.org/10.1007/978-3-540-69568-4_2}{10.1007/978-3-540-69568-4\_2}.

\bibitem{openwakeword2023}
D.~Hughes,
``OpenWakeWord: An open-source wake word detection framework,''
GitHub, 2023.
\newblock [Online]. Available: \url{https://github.com/dscripka/openWakeWord}

\bibitem{rudzicz2012torgo}
F.~Rudzicz, A.~K.~Namasivayam, and T.~Wolff,
``The TORGO database of acoustic and articulatory speech from speakers with
dysarthria,''
\textit{Language Resources and Evaluation}, vol.~46, no.~4,
pp.~523--541, 2012.
\newblock DOI: \href{https://doi.org/10.1007/s10579-011-9145-0}{10.1007/s10579-011-9145-0}.

\end{thebibliography}

% ============================================================
% GLOSSARY
% ============================================================
\newpage
\onecolumn
\printglossary[type=\acronymtype, title=List of Acronyms,
               toctitle=List of Acronyms]
\printglossary[title=Terminology, toctitle=Terminology]

\end{document}